\section{Problem Model}\label{sec:formulation}

The input is a micro-operation DAG $G=(\mathcal V,\mathcal E)$. Operation nodes
carry a pipeline, cycle count, and buffer set; cache-management nodes give each
buffer $b$ a unique allocation $a_b$, free $f_b$, size $s_b$, and cache type
$\tau(b)$. A base schedule is a topological order of all original nodes.
Evaluation uses L1${}=4096$, UB${}=1024$, L0A${}=256$, L0B${}=256$, and
L0C${}=512$.

\begin{figure*}[t]
  \centering
  \begin{tikzpicture}[
      font=\footnotesize,
      op/.style={draw=ksTikzOp!85!black, fill=ksTikzOp!13, rounded corners=2.5pt,
                 minimum height=8.6mm, minimum width=21mm, align=center,
                 line width=0.7pt},
      mgmt/.style={draw=ksTikzMg!70!black, fill=ksTikzMg!9, rounded corners=2.5pt,
                   minimum height=8.6mm, minimum width=21mm, align=center,
                   line width=0.7pt},
      dep/.style={-{Stealth[length=2.4mm,width=2mm]}, line width=0.7pt,
                  black!60},
      life/.style={{Bar[width=1.5mm]}-{Bar[width=1.5mm]}, line width=1.1pt},
      note/.style={font=\footnotesize, anchor=west, text=black}
    ]
    \definecolor{ksTikzOp}{HTML}{4C72B0}
    \definecolor{ksTikzMg}{HTML}{C44E52}
    \definecolor{ksTikzGen}{HTML}{DD8452}

    % DAG nodes
    \node[mgmt] (a0) at (0,0)    {\textsc{alloc} $b_0$\\[-1.5pt] {\scriptsize\color{black}backed}};
    \node[op]   (ci) at (3.1,0)  {\textsc{copy-in}\\[-1.5pt] $\{b_0\}$};
    \node[op]   (cp) at (6.2,0)  {\textsc{compute}\\[-1.5pt] $\{b_0,b_1\}$};
    \node[op]   (co) at (9.3,0)  {\textsc{copy-out}\\[-1.5pt] $\{b_1\}$};
    \node[mgmt] (a1) at (3.1,-1.8) {\textsc{alloc} $b_1$\\[-1.5pt] {\scriptsize\color{black}generated}};
    \node[mgmt] (f0) at (6.2,-1.8) {\textsc{free} $b_0$};
    \node[mgmt] (f1) at (9.3,-1.8) {\textsc{free} $b_1$};

    % dependency edges
    \draw[dep] (a0) -- (ci);
    \draw[dep] (ci) -- (cp);
    \draw[dep] (cp) -- (co);
    \draw[dep] (a1) -- (cp);
    \draw[dep] (cp) -- (f0);
    \draw[dep] (co) -- (f1);

    % logical lifetimes with spill-class costs
    \draw[life, ksTikzOp]  (0,-3.0)  -- (6.2,-3.0);
    \node[note, anchor=north, align=center] at (3.1,-3.12)
      {$b_0$: backed ($\kappa_{b_0}{=}1$) $\Rightarrow$ $1\times s_{b_0}$ (reload only)};
    \draw[life, ksTikzGen] (3.1,-3.72) -- (9.3,-3.72);
    \node[note, anchor=north, align=center] at (6.2,-3.84)
      {$b_1$: generated ($\kappa_{b_1}{=}2$) $\Rightarrow$ $2\times s_{b_1}$ (write-back $+$ reload)};

    % legend
    \node[mgmt, minimum width=4.5mm, minimum height=3.2mm] at (0.1,-4.60) {};
    \node[note] at (0.4,-4.60) {cache management};
    \node[op, minimum width=4.5mm, minimum height=3.2mm] at (3.15,-4.60) {};
    \node[note] at (3.45,-4.60) {operation};
    \node[note, anchor=west] at (5.25,-4.60)
      {Canonical P2 identity: $\mathrm{Tr}=\mathrm{Cl}+2\mathrm{Dt}=(\mathrm{Cl}+\mathrm{Dt})+\mathrm{Dt}=\mathrm{Vol}+\mathrm{Dt}$};
  \end{tikzpicture}

  \caption{Micro-operation DAG and benchmark spill classes. Cache-management
  nodes delimit logical lifetimes. Input-backed buffers incur one transfer per
  interruption; generated, unbacked buffers incur two.}
  \Description{A small allocation-operation-free DAG with one backed and one
  generated buffer, their lifetimes, and spill costs.}
  \label{fig:dag-example}
\end{figure*}

The evaluator assigns a \emph{static} class: if $b$ appears in any
\textsc{copy-in}, then $\kappa_b=1$ (backed); otherwise $\kappa_b=2$
(generated and unbacked). One spill episode costs
$\mathrm{tr}_b=\kappa_b s_b$. Since the input does not expose per-operation
read/write roles, this is an evaluator definition, not dynamic dirty-state
inference. Let backed and generated spill volumes be $\mathrm{Cl},\mathrm{Dt}$:
\begin{equation}\label{eq:extra-decomp}
 \mathrm{Tr}=\mathrm{Cl}+2\,\mathrm{Dt}=\mathrm{Vol}+\mathrm{Dt},
 \qquad \mathrm{Vol}=\mathrm{Cl}+\mathrm{Dt}.
\end{equation}
Thus $\mathrm{Vol}\le\mathrm{Tr}\le2\mathrm{Vol}$. At fixed volume, class mix
changes traffic by at most two, while ordering and placement can also change
$\mathrm{Vol}$; tightness and scope appear in Supplementary Sec.~2.1.

A physical residency segment occupies a contiguous interval in its cache, and
simultaneous segments may not overlap. A \textsc{spill-out}/\textsc{spill-in}
pair interrupts residency, possibly reloading at another address. Original
users and \textsc{free} are mandatory-residency events. P1 scores logical
\textsc{alloc}--\textsc{free} peaks only. P2 emits the augmented schedule,
addresses, and spill table and uses $(\mathrm{Tr},n,T)$; P3 uses
$(T,\mathrm{Tr},n)$, where $T$ is pipelined earliest completion. P1 logical
peaks are not P2 physical peaks, and P3 may accept traffic to reduce time, so
the views are never mixed. Every physical result is checked for coverage,
topology, contiguous packing, and spill dependencies. The complete evaluator
contract and paired P2--P3 tradeoff are in Supplementary Secs.~1 and~4.4.
